% !TEX root = heat_from_bit_ru.tex
\documentclass[12pt,a4paper]{article}
\usepackage{cmap}
\usepackage[margin=1in]{geometry}
\usepackage{amsmath,amssymb,amsfonts}
\usepackage[utf8]{inputenc}
\usepackage[T2A]{fontenc}
\usepackage[russian]{babel}
\usepackage[hidelinks]{hyperref}
\newcommand{\Sref}[1]{\hyperref[#1]{\S\ref{#1}}}
\usepackage{comfortaa}
\renewcommand*\familydefault{\sfdefault}
\usepackage[x11names]{xcolor}
\usepackage{indentfirst}
\usepackage{booktabs}
\usepackage{array}
\usepackage{graphicx}
\usepackage{tabularx}
\usepackage{ragged2e}
\usepackage{makecell}
\usepackage{colortbl}
\arrayrulecolor{white}

\makeatletter
\renewcommand{\thefootnote}{\arabic{footnote}}
\renewcommand{\@makefnmark}{\mbox{\textsuperscript{\textcolor{white}{\thefootnote}}}}
\renewcommand{\@makefntext}[1]{\parindent 1.8em\noindent\textsuperscript{\textcolor{white}{\thefootnote}}\ \textcolor{white}{#1}}
\renewcommand{\footnoterule}{\kern-3pt\color{white}\hrule width 0.4\textwidth height 0.4pt\kern 2.6pt}
\renewcommand{\ps@plain}{%
  \let\@oddhead\@empty
  \let\@evenhead\@empty
  \def\@oddfoot{\hfill\textcolor{white}{\thepage}\hfill}%
  \let\@evenfoot\@oddfoot
}
\makeatother

\pagestyle{plain}

\title{\textbf{Heat from Bit: термодинамическая цена объективного коллапса}}
\author{%
  Олег Понфиленок\\[0.35em]%
  {\normalsize Независимый исследователь; \href{mailto:ponfil@gmail.com}{ponfil@gmail.com}}%
}
\date{}

\begin{document}

\pagecolor{black}
\color{white}

\maketitle

\begin{abstract}
Выдвигаем гипотезу: любой механизм \emph{объективного} коллапса обязан заплатить термодинамическую цену $k\ln 2$ за уничтоженный бит когерентности \emph{в pointer-базисе акта} --- как тепло $\Delta Q = kT\ln 2$ в среде с температурой~$T$. Тариф вычисляется на состоянии, обусловленном только физически доставленным локальным pointer-регистром. Доказываем, что нетривиальный неотрицательный тариф, нулевой на pointer-диагональных состояниях, не может быть аффинным функционалом одной локальной матрицы плотности. Проверка: калориметрия нагрева при поглощении геральдированного однофотонного кубита --- $|{+}\rangle$ против контроля $I/2$.
\end{abstract}

\noindent\textbf{Ключевые слова:} объективный~коллапс; pointer-состояние; квантовый~дарвинизм; принцип~Ландауэра; тепло~из~бита; инкассация~энтропии

%%% ===================================================================
\section{Введение}\label{sec:intro}
%%% ===================================================================

Идея \emph{объективного} коллапса~--- сделать редукцию состояния физическим процессом, а не эпистемическим постулатом измерения~--- образует отдельную исследовательскую программу на фоне стандартной квантовой механики, где кажущийся коллапс объясняют унитарной декогеренцией и корреляциями с окружением. Ранняя линия Каройхази~\cite{karolyhazy1966} связала гибель суперпозиций с неопределённостью пространства-времени; GRW и CSL~\cite{grw1986,pearle1989} ввели спонтанную локализацию; модели Диоши--Пенроуза~\cite{diosi1989,penrose1996}~--- гравитационную неустойчивость макроскопических суперпозиций; в других моделях: пересечение порога скремблирования~\cite{ponfilenok2026threshold}, ветвезависимая дефазировка собственного времени~\cite{ponfilenok2026branch}. Обзор и статус~--- у Bassi et al.~\cite{bassi2013}: направление активно развивается и экспериментально ограничивается, но остаётся маргинальным относительно ортодоксальной квантовой теории; решающего подтверждения или окончательного опровержения пока нет. Настоящая работа не выбирает критерий «когда», а ставит вопрос о термодинамической цене любого такого акта.

Объективный коллапс вводит фундаментально необратимый элемент динамики, отличающий выделенное направление времени от обратимой унитарной эволюции. Однако такая необратимость сама по себе не эквивалентна нарушению непрерывной симметрии временных трансляций, с которой теорема Нётер связывает сохранение энергии. Поэтому энергетический след коллапса зависит от конкретной неунитарной динамики; настоящая работа выдвигает гипотезу о его термодинамической природе (\Sref{sec:origin}).

Принцип Ландауэра~\cite{landauer1961} связывает уничтожение \emph{классической} информации с минимальной диссипацией $Q \ge kT\ln 2$. Плата сидит на логически необратимом сжатии фазового объёма (многие$\to$один)~--- следствие второго начала~\cite{bennett1982,sagawa2008}. При декогеренции среда выбирает \emph{pointer-базис}~--- набор устойчивых исходов $|i\rangle$, размножаемых в окружении (einselection; квантовый дарвинизм~\cite{zurek2009}). Pointer-базис задаёт множество \emph{классически допустимых} (incoherent) состояний~--- диагональных смесей по $|i\rangle$. \emph{Относительная энтропия когерентности}~\cite{baumgratz2014} $C_r(\rho)=\min_{\sigma\in\mathcal{I}} S(\rho\|\sigma)$ измеряет, насколько $\rho$ не сводится к такой смеси; для pointer-базиса минимум достигается на $\rho_{\mathrm{diag}}$ и $C_r(\rho)=S(\rho_{\mathrm{diag}})-S(\rho)$. У чистого $|{+}\rangle_{\mathrm{ptr}}$ $S(\rho)=0$, но $C_r=k\ln 2$~--- полная which-path когерентность; у $I/2_{\mathrm{ptr}}$ $S(\rho)=k\ln 2$, но $C_r=0$~--- та же pointer-статистика без фаз, стирать нечего. При \emph{унитарной} декогеренции $|{+}\rangle \to I/2$ $S(\rho_S)$ подсистемы растёт до $k\ln 2$, но энтропия целого нулева: фаза уходит в корреляции и восстановима (эхо Лошмидта~\cite{gorin2006,goussev2012}). При \emph{объективном} коллапсе pointer-исход зафиксирован необратимо, $H=S_{\mathrm{coh}}=C_r$ в pointer-базисе~--- термодинамически активная величина; по гипотезе настоящей работы её скачок материализуется как тепло. Hormoz~\cite{hormoz2013} показывает необходимость энтропийной платы за коллапс в ratchet-модели; Heat from Bit квантифицирует плату для объективной фиксации pointer-исхода без явного демона.

Уилер сформулировал «It from Bit»~\cite{wheeler1989}: физическая реальность («it») возникает из актов различения pointer-исходов. \emph{Heat from Bit}~--- та же параллель в термодинамике: $C_r(\rho)$ в pointer-базисе~--- «bit»; её необратимый скачок~--- «it», материализующийся как тепло.

Предлагаем экспериментальную проверку (§\,\ref{sec:experiment}): калориметрия \emph{геральдированного} однофотонного кубита в интерферометре М--Z (heralded MZ) --- $|{+}\rangle_{\mathrm{ptr}}$ даёт $kT\ln 2$ на акт, классическая смесь $I/2_{\mathrm{ptr}}$~--- нет.

%%% ===================================================================
\section{Интуиция}\label{sec:intuition}
%%% ===================================================================

Идея Heat from Bit выросла из простого мысленного эксперимента, в котором \emph{один} фотон поглощается анизотропным поглотителем в базисе поляризации, задающим pointer-базис акта. Базис~--- $\{|0\rangle,|1\rangle\}$ линейной поляризации ($|0\rangle$~--- горизонталь, $|1\rangle$~--- вертикаль): какой из этих исходов необратимо зафиксируется при поглощении~--- то есть при детекции?

\textbf{Случай~A (определённая поляризация).} Однофотонное состояние $|0\rangle$~--- собственное состояние pointer-базиса, совпадающее с осью поляризатора. Когерентность $S_{\mathrm{coh}}=0$. Фотон поглощается, поглотитель получает энергию $E_\gamma$~--- обычное первое начало; альтернативный pointer-исход $|1\rangle$ в предактовом состоянии отсутствует, стирать нечего. При акте объективного коллапса ничего не меняется.

\textbf{Случай~B (суперпозиция pointer-исходов).} Тот же однофотонный протокол, но фотон подготовлен в
\begin{equation}
|{+}\rangle_{\mathrm{ptr}} = \frac{|0\rangle + |1\rangle}{\sqrt{2}}
\end{equation}
в pointer-базисе поляризатора: $S_{\mathrm{coh}}=\ln 2$. Поглощение передаёт ту же $E_\gamma$; различие~--- в предактовом состоянии, а не в энергии фотона.

\emph{Шаг~1 (поглощение).} Независимо от подготовки акт поглощения передаёт поглотителю $E_\gamma$; объективного коллапса на этом этапе ещё нет.

\emph{Шаг~2 (декогеренция).} При \emph{унитарной} эволюции когерентность в pointer-базисе $S_{\mathrm{coh}}$ уходит во внутренние степени свободы поглотителя и среды~--- в запутанность, восстановимую в принципе (эхо, обратимый цикл). Среда уже взаимодействует с фотоном и переносит сведения о pointer-исходе (здесь~--- о поляризации) во \emph{внутреннюю память} окружения. Запутанность поглотителя и среды с $|0\rangle$ и $|1\rangle$ \emph{и есть} эта память~--- альтернатива не уничтожена, а закодирована в корреляциях; при унитарной эволюции она сохраняется, стирания нет, глобальная энтропия не растёт.

\emph{Шаг~3 (объективный коллапс).} Объективный коллапс можно понимать как \emph{очистку} внутренней памяти: степени свободы, нёсшие когерентность и запутанность, исчезают необратимо~--- процедура стирания регистра, а не обратимый перенос фазы. Pointer-исход зафиксирован; ``внутреннего'' квантового состояния, хранившего альтернативу $|0\rangle$ vs $|1\rangle$, больше нет, $S_{\mathrm{coh}}$ исчезает. Это~--- акт стирания одного бита pointer-исхода: из суперпозиции $|{+}\rangle$ необратимо остаётся один исход $|0\rangle$ или $|1\rangle$. Мы выдвигаем гипотезу, что фундаментальное уничтожение этой информации подчиняется термодинамическому ограничению, аналогичному принципу Ландауэра~\cite{landauer1961}: среда акта~--- резервуар, где жили pointer-моды и их корреляции с поглотителем,~--- должна получить дополнительное тепло
\begin{equation}
\Delta Q = kT\ln 2
\end{equation}
\emph{на каждый} объективно зафиксированный однофотонный акт ($|{+}\rangle_{\mathrm{ptr}}\to|0\rangle$ или $|1\rangle$), сверх обычного нагрева на $E_\gamma$. Избыток случая~B над~A при равной $E_\gamma$~--- предсказание Heat from Bit.

\textbf{Случай~C (классическая смесь).} Тот же протокол, но фотон подготовлен как $I/2$ в pointer-базисе:
\begin{equation}
I/2 = \frac{|0\rangle\langle 0| + |1\rangle\langle 1|}{2}
\end{equation}
--- диагональная смесь $|0\rangle$ и $|1\rangle$ без фазовых связей: $S_{\mathrm{coh}}=0$. Это не суперпозиция pointer-исходов, а статистическая запись: исход уже ``измерен'' на уровне предактового состояния, остаётся лишь классическая неопределённость \emph{какой} из них реализовался. Среда в акте поглощения не \emph{стирает} альтернативный бит, а лишь \emph{копирует} уже имеющуюся классическую запись; копирование классической информации энергетически бесплатно. Ни декогеренции, ни объективного коллапса в акте нет.

%%% ===================================================================
\section{Pointer-базис и энтропия когерентности}\label{sec:pointer}
%%% ===================================================================

При взаимодействии с окружением система декогерирует в \emph{pointer-базис} $\{|i\rangle\}$~--- устойчивых исходов, размножаемых в среде (einselection; квантовый дарвинизм~\cite{zurek2009}). Классическая запись~--- конкретный pointer-исход $|i\rangle$. $|{+}\rangle$ и $I/2$ совпадают по $\rho_{\mathrm{ptr}}$ и $\langle\hat{H}\rangle$; болометр видит лишь диагональ в pointer-базисе. Тариф~--- по когерентности в pointer-базисе \emph{акта}, а не по $S(\rho)$ абстрактно.


\textbf{Когерентность в pointer-базисе.} В ресурсной теории когерентности~\cite{baumgratz2014,streltsov2017}:
\begin{equation}
S_{\mathrm{coh}}(\rho) = S(\rho_{\mathrm{ptr}}) - S(\rho),
\end{equation}
где $\rho_{\mathrm{ptr}}$~--- диагональ $\rho$ в pointer-базисе акта. Для $|{+}\rangle$ $S_{\mathrm{coh}}=\ln 2$, для $I/2$~--- ноль; это \emph{relative entropy of coherence} $C_r(\rho)$~\cite{baumgratz2014,streltsov2017} относительно того же базиса. Kammerlander и Anders~\cite{kammerlander2016} показывают, что термодинамика \emph{проекции} в базис (дефазировка, неселективное измерение) нетривиальна только при когерентности в этом базисе; для диагональных смесей $\Delta S=0$. В оптимальном процессе для $|{+}\rangle$ из такой проекции можно извлечь $\langle W\rangle=kT\ln 2$ работы. Heat from Bit переносит эту различимость на pointer-базис акта и постулирует связанную с $S_{\mathrm{coh}}$ диссипацию при объективной фиксации исхода.

\textbf{Тарифная энтропия $H$.} Обозначим $H$~--- \emph{тарифную энтропию} предактового состояния в pointer-базисе \emph{акта}; $S(\rho)$, $S_{\mathrm{coh}}$ и $H$~--- в \emph{натах}:
\begin{equation}
H(\rho) = S_{\mathrm{coh}}(\rho).
\end{equation}
$H$~--- не полная $S(\rho)$ и не классическая $-\sum_i p_i\ln p_i$, а именно когерентность в pointer-базисе, уничтожаемая при необратимой фиксации исхода, задаёт термодинамическую плату.
Диагональ $\rho_{\mathrm{ptr}}$ ($I/2$): $H=0$; чистая суперпозиция $|{+}\rangle_{\mathrm{ptr}}$: $H=\ln 2$.

\textbf{Число бит}:
\begin{equation}
N_{\mathrm{bit}}(\rho) = H(\rho) / \ln 2.
\end{equation}

\textbf{Когерентные состояния поля} ($\hat a|\alpha\rangle=\alpha|\alpha\rangle$) --- собственные состояния оператора $\hat a$~\cite{glauber1963}; они задают классическую амплитуду моды и CW-лазерный луч. Термин ``когерентный'' здесь не означает $S_{\mathrm{coh}}$ в pointer-базисе: в $|\alpha\rangle$ нет суперпозиции pointer-исходов, только пуассоновская статистика, $H=0$; у ``фотона'' когерентного луча нет отдельного предактового $\rho$. ``Угол $45^\circ$'' у CW-лазера не воспроизводит случай~B из \Sref{sec:intuition}:
\begin{equation}
|\alpha\rangle_{45^\circ} = |\alpha/\sqrt{2}\rangle_{0}\,|\mathrm{i}\alpha/\sqrt{2}\rangle_{1},
\end{equation}
а не поток $|{+}\rangle_{\mathrm{ptr}}$. Для чистого предактового $\rho=|\psi\rangle\langle\psi|$, $|\psi\rangle=\sum_i c_i|i\rangle$, имеем $p_i=|c_i|^2$ и $H=-\sum_i p_i\ln p_i$ (тогда $S(\rho)=0$). В общем случае $H=S(\rho_{\mathrm{ptr}})-S(\rho)$; для смешанного $I/2$ это $H=0$, хотя $-\sum_i p_i\ln p_i=\ln 2$.

%%% ===================================================================
\section{Термодинамика объективного коллапса}\label{sec:thermo}
%%% ===================================================================

Тариф $H$ (\Sref{sec:pointer})~--- определение ($H=S_{\mathrm{coh}}$ в pointer-базисе акта), а не предположение модели.

При \emph{унитарной} дефазировке $S(\rho_{\mathrm{ptr}})$ не меняется, рост $S(\rho)$ подсистемы~--- убыль $S_{\mathrm{coh}}$ (\Sref{sec:pointer}); изолированная подсистема при унитарном $\pi/2$ не меняет $S(\rho)$ (чистое остаётся чистым), а обратимый цикл $\pi/2\to-\pi/2$ восстанавливает состояние; тариф $H$ не активируется. Тепло \emph{не} следует из роста $S(\rho)$ подсистемы и не сопровождает унитарное «шифрование» фазы pointer-базиса в корреляциях с окружением.

\textbf{Постулат (Heat from Bit).} Проверяемая гипотеза работы (\Sref{sec:experiment}): при \emph{объективном} коллапсе с $H>0$ скачок тарифа отождествляется с термодинамической энтропией резервуара $S_{\mathrm{th}}=kH$ и теплом $\Delta Q=T\Delta S_{\mathrm{th}}=kT\Delta H$ в среде при температуре~$T$; pointer-исход зафиксирован необратимо. Среднее по актам
\begin{equation}
\langle Q\rangle = kT\,H = kT\,C_r(\rho).
\label{eq:tariff_mean}
\end{equation}
При данном $H$ распределение $Q_i$ из этого \emph{не} следует (например, $Q_i \equiv kT H$ даёт то же $\langle Q\rangle$). Для \emph{чистого} предактового $|\psi\rangle=\sum_i c_i|i\rangle$ ($S(\rho)=0$, $p_i=|c_i|^2$) модель \emph{дополнительно} принимает потраекторный тариф
\begin{equation}
\boxed{Q_i = kT\ln(1/p_i)}
\end{equation}
--- информационная ``неожиданность'' исхода~$i$ (согласованная со средним через $H=-\sum p_i\ln p_i$): редкий исход ``стирает'' больше альтернативной вероятностной массы. Тогда $\Delta S_{\mathrm{th},i}=k\ln(1/p_i)$,
\begin{equation}
\langle Q\rangle = \sum_i p_i\, Q_i = kT H, \qquad \mathrm{Var}(Q) = (kT)^2\,\mathrm{Var}(\ln p_i).
\end{equation}
Для смешанного состояния, диагонального в pointer-базисе ($H=0$, например $I/2$), поактный тариф не активируется: $\langle Q\rangle=0$; частично когерентная смесь в общем случае имеет $H>0$. При равных $p_i=1/n$ все исходы чистого состояния эквивалентны: $Q_i = kT\ln n$; для бинарного $|{+}\rangle$ ($p_i=\tfrac12$) это $\Delta Q = kT\ln 2$, $\Delta H=\ln 2$ и $\Delta S_{\mathrm{th}} = k\ln 2$ ($\Delta N_{\mathrm{bit}}=1$). Пример $p_0=0{,}9$, $p_1=0{,}1$: акт с исходом~0 выделяет $kT\ln(10/9)\approx 0{,}11\,kT$, с исходом~1~--- $kT\ln 10\approx 2{,}3\,kT$; среднее $\approx 0{,}33\,kT = kT\,h_2(0{,}9)$.

Тариф по исходам согласован с правилом Борна и стохастической термодинамикой~\cite{sagawa2008}. При равных амплитудах ($|{+}\rangle$) $Q_i$ одинаково, флуктуаций нет; различие версий~--- лишь при несимметричной суперпозиции и поактной калориметрии (пока недостижимой; \Sref{sec:experiment}).

Распространение уже существующей записи (CNOT, избыточность в смысле квантового дарвинизма~\cite{zurek2009}) нового бита не порождает; унитарная запутанность (GHZ, кот-состояния) тариф не умножает. Унитарное приготовление $|{+}\rangle_{\mathrm{ptr}}$ ($\pi/2$, which-path-кодирование $|x_0\rangle+|x_1\rangle$ и т.п.)~--- обратимые операции; платы нет.

На серии объективных актов (среднее по исходам):
\begin{equation}
dH = \ln 2\,dN_{\mathrm{bit}}, \qquad \langle\delta Q\rangle = kT\,dH = \mu_{\mathrm{bit}}\,dN_{\mathrm{bit}}, \qquad \mu_{\mathrm{bit}} = kT\ln 2.
\end{equation}
Для чистого предактового состояния один акт с исходом~$i$: $\Delta Q_i = kT\ln(1/p_i)$; $\langle\Delta Q\rangle = kT\,\Delta H = kT\ln 2\,\Delta N_{\mathrm{bit}}$. Платёж~--- не $\langle\hat{H}\rangle$ гамильтониана, а термодинамическая цена фиксации pointer-исхода; $\delta Q$ рассеивается в резервуаре, в среднем~--- по тарифу $\mu_{\mathrm{bit}}$.

\subsection{Состояние, на котором вычисляется тариф}\label{sec:conditioning}

Формула~\eqref{eq:tariff_mean} неполна без правила, какое именно локальное состояние входит в $C_r$. Различим три возможности. (i)~\emph{Безусловное локальное правило}: $Q=kT C_r(\rho_S)$ для приведённой матрицы системы акта. Оно причинно и не различает разные ансамблевые разложения одной $\rho_S$, но фазовый scrambler после усреднения становится тем же $I/2$, что и классический контроль. (ii)~\emph{Локально обусловленное правило}: если в причинное состояние области акта физически вошёл объективный классический pointer-регистр $R$ с записью $r$, используется условное состояние $\rho_{S|r}$,
\begin{equation}
\langle Q\rangle_R=kT\sum_r p_r\,C_r(\rho_{S|r}).
\label{eq:local_conditioned_tariff}
\end{equation}
(iii)~\emph{Глобально обусловленное правило}: выбирается удалённая онтическая ветвь даже без локально доставленной записи; эта более сильная версия может различать удалённо приготовленные разложения одной $\rho_S$.

В настоящей работе принимается правило~(ii). Здесь ``запись'' означает физический локальный регистр, причинно доставленный к акту, а не знание наблюдателя и не информацию, лишь существующую где-то в прошлом световом конусе. Если $R$ сам остаётся в суперпозиции или только запутан с $S$, условный классический индекс $r$ не вводится: произвольное ансамблевое разложение $\rho_S$ выбирать нельзя. Правило~\eqref{eq:local_conditioned_tariff} поэтому требует внешнего критерия, когда $R$ стал объективным pointer-регистром; Heat from Bit задаёт цену такого акта, но не универсальный критерий ``когда''.

\paragraph{Теорема о линейном тарифе.}
Пусть сектор одного акта конечномерен, $\{|i\rangle\}$~--- полный pointer-базис, а средний тариф аффинен по локальной матрице состояния. Тогда $Q(\rho)=\operatorname{Tr}(X\rho)$ для некоторого эрмитова $X$ (константа включается в компоненту, пропорциональную единице). Если
\begin{equation}
Q(|i\rangle\langle i|)=0\quad\forall i,
\qquad
Q(\rho)\ge0\quad\forall\rho,
\label{eq:linear_tariff_conditions}
\end{equation}
то $X=0$ и $Q\equiv0$. Действительно, второе условие означает $X\ge0$, первое даёт $X_{ii}=0$, а положительность влечёт $|X_{ij}|^2\le X_{ii}X_{jj}=0$. Таким образом, \emph{нетривиальный} положительный тариф, нулевой на всех pointer-диагональных состояниях, не может быть аффинным функционалом одной локальной $\rho$. Это теорема о тривиальности линейной ветви, а не эквивалентность ``нелинейность $\Leftrightarrow$ сигнализация''.

Аффинность локального наблюдаемого достаточна, чтобы разные разложения одной $\rho_S$ не давали сигнал через этот канал; обратный вывод неверен. В частности, нелинейный $kT C_r(\rho_S)$ также не сигнализирует: для половины пары Белла $\rho_B=I/2$ при любом удалённом базисе и $C_r(I/2)=0$. Сигнал появляется лишь при дополнительном глобальном обусловливании~(iii), а не из нелинейности самой по себе~\cite{gisin1989,gisin1990,polchinski1991,hjw1993}.

%%% ===================================================================
\section{Происхождение энергии тепла}\label{sec:origin}
%%% ===================================================================

В конкретных моделях объективного коллапса неунитарная динамика сама может задавать энергию, выделяющуюся при локализации: в CSL и GRW~\cite{grw1986,pearle1989,bassi2013} оператор коллапса повышает среднюю энергию системы. Тогда тарифное тепло $\delta Q = kT\Delta H$ (\Sref{sec:thermo}) может, в~принципе, идти из этого резерва. В других моделях такой резерв может не выделяться явно или оказываться пренебрежимо малым.

Мы рассматриваем длительность акта коллапса $\tau_{\mathrm{col}}$~--- неунитарного сжатия суперпозиции к одному pointer-исходу. При декогеренции система запутывается со степенями свободы среды акта~--- ``приёмниками'' pointer-информации; резервуар при температуре~$T$ поглотителя наполнен модами на масштабе~$kT$, характерное время тепловых колебаний~--- $\tau_{\mathrm{th}} \equiv \hbar/kT$. Во многих теориях объективного коллапса момент локализации зависит от числа частиц среды, с которыми установилась декогеренция, или от эффективной массы суперпозиции~\cite{grw1986,bassi2013,zurek2009,ponfilenok2026threshold,ponfilenok2026branch}. Эта масса нарастает по мере присоединения новых ``приёмников'', подключающихся за время порядка~$\tau_{\mathrm{th}}$. Время $\tau_{\mathrm{col}}$ задаёт конкретная модель коллапса, однако, эвристически, при критерии, растущем с числом приёмников, неунитарное сжатие должно завершиться за $\tau_{\mathrm{col}} \lesssim \tau_{\mathrm{th}}$~--- пока к суперпозиции не добавилась ещё одна степень свободы среды и критерий не сдвинулся.

Унитарная эволюция с не зависящим от времени гамильтонианом сохраняет $\langle\hat{H}\rangle$ точно. Неунитарный акт длительности $\tau_{\mathrm{col}}$ этим механизмом не связан и \emph{может} менять энергию. Соотношение Мандельштама--Тамма~\cite{mandelstam1945} задаёт не величину $\Delta\langle\hat{H}\rangle$ (это был бы нагрев в духе GRW), а лишь масштаб допуска, на котором о сохранении для данного акта ещё имеет смысл говорить: $\Delta E\sim\hbar/(2\tau_{\mathrm{col}})$. Чтобы тарифный платёж $kT\ln 2$ не превосходил этого окна, свяжем $\tau_{\mathrm{col}}$ с $\delta Q = kT\ln 2$:
\begin{equation}
\tau_{\mathrm{col}} \lesssim \frac{\hbar}{2kT\ln 2} \approx \frac{\tau_{\mathrm{th}}}{2\ln 2},
\label{eq:tau_col}
\end{equation}
Если платёж сидит внутри окна, локальный дисбаланс не требует отдельного резерва; после акта остаток может рассеяться в тепловом резервуаре. Это верхняя граница согласованности, а не расчёт нагрева. Отдельный акт микроскопически не аудируется; по серии из $N$ актов наблюдаем суммарный избыток $\Delta Q_{\mathrm{tot}} = N\,kT\Delta H$ (\Sref{sec:experiment}).

%%% ===================================================================
\section{Экспериментальное предложение}\label{sec:experiment}
%%% ===================================================================

Предложенный тест~--- дифференциальный: постулируемый член \emph{добавляет} тепло к обычному нагреву ($Q_{\mathrm{out}} = Q_{\mathrm{std}} + \mu_{\mathrm{bit}}\,\Delta N_{\mathrm{bit}}$ при $\Delta N_{\mathrm{bit}} > 0$), поэтому измеряется систематический избыток над контролем с $H=0$. Здесь $Q_{\mathrm{std}}$~--- тепло без тарифа (поглощение $W_{\mathrm{in}}$, обычная $Q_{\mathrm{dissip}}$); $\Delta N_{\mathrm{bit}}$~--- число pointer-битов, объективно стёртых в актах (\Sref{sec:pointer}): один акт из $|{+}\rangle_{\mathrm{ptr}}$ даёт $\Delta N_{\mathrm{bit}}=1$, из $I/2_{\mathrm{ptr}}$~--- ноль. Heralded MZ~--- $\sim 10^{10}$ при $R=10^{8}$~с$^{-1}$ и $\tau\sim 10^{2}$~с; все проверяют $\langle Q\rangle = kT H$ (\Sref{sec:thermo}), не различая флуктуации $Q_i$ по исходам. Тест корреляции $Q$ с исходом требует несимметричной суперпозиции и поактной калориметрии и \emph{пока} недостижим: $\Delta Q_i \sim 10^{-2}$~эВ на herald при $T=300$~K сопоставимо с шумом однособытийного болометра, а heralded MZ нацелен на $\langle Q\rangle$ через lock-in и накопление $\sim 10^{10}$ актов.


Обычный CW-лазер после 50/50~--- произведение когерентных мод $|\alpha/\sqrt{2}\rangle_L|\mathrm{i}\alpha/\sqrt{2}\rangle_R$, а не поток $|{+}\rangle_{\mathrm{ptr}}$-кубитов; то же для поляризации $|\alpha\rangle_{45^\circ}$ (\Sref{sec:pointer}). Здесь~--- чистый поток кубитов: SPDC + SNSPD-геральд $\to$ $|1\rangle_{\mathrm{Fock}}$ $\to$ Mach--Zehnder $\to$ две pointer-зоны поглощения.

\paragraph{Pointer-базис и состояния.}
Pointer-базис акта~--- \emph{локализация поглощения}: $|x_0\rangle$, $|x_1\rangle$~--- две разнесённые зоны на одном микроострове (левое и правое плечо М--Z). Однофотонное path-состояние после первого делителя М--Z:
\begin{equation}
|\psi\rangle = \frac{|1,0\rangle + e^{\mathrm{i}\phi}|0,1\rangle}{\sqrt{2}}
\quad\Leftrightarrow\quad
|{+}\rangle_{\mathrm{ptr}}:\;
H = S_{\mathrm{coh}} = C_r = \ln 2.
\end{equation}
Приготовление унитарно ($S(\rho)=0$); тариф активируется только при объективном коллапсе (\Sref{sec:thermo}), не при делении или фазовом повороте. Контроль~--- \emph{классическая} подготовка path без когерентности (не phase-scrambler):
\begin{equation}
\rho_{\mathrm{mix}} = \tfrac{1}{2}|1,0\rangle\langle 1,0| + \tfrac{1}{2}|0,1\rangle\langle 0,1|
\quad\Leftrightarrow\quad
I/2_{\mathrm{ptr}}:\;
H = 0;
\end{equation}
диагонали $\rho_{\mathrm{ptr}}$ совпадают с $|{+}\rangle_{\mathrm{ptr}}$, но $C_r(\rho_{\mathrm{mix}})=0$. Контроль реализуют \emph{до} интерферометра: оптический переключатель по объективному локальному биту $r$ (физически $|1,0\rangle$ или $|0,1\rangle$ на каждый herald). По принятому локально обусловленному правилу~\eqref{eq:local_conditioned_tariff} условное состояние уже является pointer-состоянием, поэтому $H=0$.

Быстрая рандомизация фазы не эквивалентна этому контролю, если значение $\phi$ объективно записано в локальном регистре модулятора и причинно входит в подготовку акта: условно на запись состояние остаётся чистым $|\psi(\phi)\rangle$ и $H=\ln2$. Если же контроллер фазы сам остаётся квантовым, его запись стёрта до акта или физически не входит в причинное состояние поглотителя, условливание запрещено и используется безусловное $I/2$ с $H=0$. Тем самым различие задаётся физическим регистром подготовки, а не выбором наблюдателя между эквивалентными разложениями одной $\rho$ (\Sref{sec:conditioning}).

\textbf{Почему which-path, а не поляризация.} \Sref{sec:pointer} показывает тот же тариф на однофотонном $|{+}\rangle_{\mathrm{ptr}}$ в базисе поляризации; heralded MZ реализует эквивалентный $|{+}\rangle_{\mathrm{ptr}}$ в базисе локализации $|x_i\rangle$~--- удобнее для event-triggered калориметрии и контроля $I/2_{\mathrm{ptr}}$.

\paragraph{Схема.}
\begin{itemize}
\item \textbf{Источник}: SPDC в монолитном или внешнем резонаторе (cavity SPDC) либо несколько параллельных пар SPDC на общий геральд; $\lambda=1550$~нм, SNSPD на idler; сигнальный $|1\rangle_{\mathrm{Fock}}$ $\to$ MZ (50/50, $\phi$ фиксирована). Проектная частота геральдов $R=10^{8}$~с$^{-1}$ (на $\sim 1$--$2$ порядка выше типичной волноводной SPDC за счёт резонатора и/или мультиплекса).
\item \textbf{Поглотитель}: подвешенный микроостров (SiN/graphene/Ge--InGaAs), две pointer-зоны на одном острове; $T=300$~K. Тепловая связь с подложкой $G\sim 1$~мкВт/К (не мВт/К: иначе фононный предел $\mathrm{NEP}=\sqrt{4k_BT^2G}$ делает lock-in неосуществимым). Квантовая эффективность тракта ``геральд $\to$ поглощение в pointer-зоне'' $\eta=0{,}85$ (SNSPD, AR-покрытия, фокус на зону).
\item \textbf{Синхронизация}: импульс геральда $\to$ окно интеграции болометра $\ll T_2^*$ path; lock-in при переключении $|{+}\rangle_{\mathrm{ptr}}\leftrightarrow I/2_{\mathrm{ptr}}$ на $f_{\mathrm{mod}}\sim 1$~кГц.
\end{itemize}

\paragraph{Сигнал.}
На \emph{каждый} объективно зафиксированный heralded-акт в канале $|{+}\rangle_{\mathrm{ptr}}$:
$\Delta Q_{\mathrm{bit}} = kT\ln 2 \approx 1{,}8\times 10^{-2}$~эВ ($\sim 2{,}2\%$ от $E_\gamma$ при $\lambda=1550$~нм); в $I/2_{\mathrm{ptr}}$~--- ноль. При проектных $R=10^{8}$~с$^{-1}$ и $\eta=0{,}85$ средний избыток мощности
\begin{equation}
\Delta P = \eta\, R\, kT\ln 2
\approx 2{,}4\times 10^{-13}~\text{Вт}
\quad (T=300~\text{K}).
\label{eq:signal_power}
\end{equation}
Доля $\Delta Q_{\mathrm{bit}}/E_\gamma$ постоянна; проверка~--- по наклону $\Delta P/R$ и по event-triggered калориметрии. При $R=10^{8}$~с$^{-1}$ примесь двухпарных SPDC-событий ($g^{(2)}(0)>0$) разбавляет дифференциал $|{+}\rangle_{\mathrm{ptr}}-I/2_{\mathrm{ptr}}$; требуется $g^{(2)}(0)\ll 1$ на herald или низкий duty cycle источника.

\paragraph{Шумовой бюджет.}
Фононный предел островка: $\mathrm{NEP}=\sqrt{4k_BT^2G}$. При проектной $G\sim 1$~мкВт/К $\mathrm{NEP}\approx 2{,}2\times 10^{-12}\;\textrm{Вт}/\sqrt{\textrm{Гц}}$. Время интегрирования lock-in до $\mathrm{SNR}\sim 1$ по мощности: $\tau\approx(\mathrm{NEP}/\Delta P)^2\sim 10^{2}$~с при $\Delta P$ выше.

\paragraph{Каналы (равный $R$, равная $\eta$).}
\textbf{Канал~$|{+}\rangle_{\mathrm{ptr}}$:} MZ с фиксированной $\phi$; $\langle Q\rangle = kT\ln 2$ на heralded-акт.

\textbf{Канал~$I/2_{\mathrm{ptr}}$:} тот же $R$; перед MZ классический which-path (переключатель L/R); $\langle Q\rangle = 0$.

Дифференциал $|{+}\rangle_{\mathrm{ptr}} - I/2_{\mathrm{ptr}}$ изолирует плату за один pointer-бит. При $R=10^{8}$~с$^{-1}$ за $\tau\sim 10^{2}$~с накапливается $\sim 10^{10}$ актов и суммарный избыток $\Delta Q_{\mathrm{tot}}\sim 10^{-11}$~Дж.

\paragraph{Что фиксируем.}
\begin{enumerate}
\item Event-triggered избыток энергии $\Delta Q_{\mathrm{bit}} \approx kT\ln 2$ на herald в канале $|{+}\rangle_{\mathrm{ptr}}$; ноль в $I/2_{\mathrm{ptr}}$.
\item Линейность суммарного $\Delta Q$ по $R$ и по $T$; наклон $\Delta P/R \propto \eta\,kT\ln 2$.
\item Lock-in-пик на $f_{\mathrm{mod}}$; отсутствие пика при постоянном $I/2_{\mathrm{ptr}}$.
\end{enumerate}
\textbf{Исходы.} Подтверждение: устойчивый дифференциал $|{+}\rangle_{\mathrm{ptr}} - I/2_{\mathrm{ptr}}$ при $R=10^{8}$~с$^{-1}$. Нуль при SNR $\gg 1$~$\Rightarrow$ bound на тариф $kT\ln 2$.

\paragraph{Причинный тест с задержанной записью.}
Правило~\eqref{eq:local_conditioned_tariff} допускает дополнительную проверку без сверхсветовой связи. Алиса измеряет свою половину пары Белла в $Z$- или $X$-базисе и передаёт Бобу классическое сообщение, содержащее \emph{и базис, и исход}. Если акт Боба происходит до физического прихода записи, локально доступно лишь $\rho_B=I/2$ и тариф равен нулю. После доставки одинакового по размеру регистра условные состояния равны $|0\rangle,|1\rangle$ для $Z$ ($H=0$) или $|+\rangle,|-\rangle$ для $X$ ($H=\ln2$). Дифференциал ``акт до/после записи'' проверяет зависимость от причинно доступной side-information; обычное тепло доставки и памяти должно сокращаться сравнением сообщений одинаковой энергии. Тот же тест можно реализовать локально коррелированным регистром, поэтому он проверяет правило обусловливания, а не квантовую нелокальность как таковую.

\paragraph{Трансмон как перспективное направление.}
Дисперсивное считывание трансмона при милли-кельвиновых температурах интересно тем, что позволяет независимо менять температуру подложки и число фотонов считывания. Однако оценивать чувствительность простой экстраполяцией $G\propto T^3$ от комнатной температуры нельзя. Квант теплопроводности $G_Q=\pi^2k_B^2T/(3h)$ является максимальным вкладом идеально передающего канала, а не универсальным нижним пределом; реальный бюджет должен включать $G$, теплоёмкость и полосу конкретного болометра, $\chi$, $\kappa$, $n_{\mathrm{crit}}$, форму импульса, фотонный shot noise и обычное состояние-зависимое тепло. Поэтому здесь трансмон фиксируется как направление для отдельного расчёта, без заявленной таблицы времени накопления~\cite{schwab2000,gambetta2006}.

%%% ===================================================================
\section{Обсуждение}\label{sec:discussion}
%%% ===================================================================

Модель Heat from Bit постулирует единственную термодинамическую плату~--- за объективный коллапс ($Q_i = kT\ln(1/p_i)$ для чистых когерентных состояний, $\langle Q\rangle = kT H$; \Sref{sec:thermo}); унитарность бесплатна. Существование таких актов принимается как предпосылка (введение); конкретный критерий «когда» статья не фиксирует. Существенная часть постулата~--- выбор состояния: рабочая версия использует локальное условливание только по физически доставленному объективному pointer-регистру (\Sref{sec:conditioning}). Потраекторный $Q_i$ остаётся самостоятельным дополнительным постулатом, а не следствием теоремы о линейном тарифе. Основная экспериментальная проверка~--- heralded MZ (\Sref{sec:experiment}).

\paragraph{Тепловой сектор и детальный баланс.}
Тепловое состояние допускает положительное $P$-представление как смесь когерентных $|\alpha\rangle$~\cite{glauber1963}, но для нелинейного тарифа это математическое разложение нельзя автоматически считать физически реализованным ансамблем. По безусловному локальному правилу тепловая $\rho$ диагональна в подходящем pointer-базисе и даёт $H=0$. По принятому условному правилу необходимо дополнительно указать, существует ли объективный локальный регистр источника и какие условные $\rho_{S|r}$ он физически готовит. Произвольный выбор удобного $P$-разложения запрещён; для квантового или удалённого регистра условливание не выполняется.

Это открытое условие совместимости с детальным балансом, а не уже доказанное освобождение равновесного света. Если каждый поглощённый солнечный фотон с характерной энергией $2$~эВ платил бы хотя бы $kT\ln2=0{,}0179$~эВ при 300~K, добавка составляла бы около $0{,}9\%$: для поглощённого потока порядка $240$~Вт/м$^2$ это $\sim2{,}1$~Вт/м$^2$, уже выше измеренного дисбаланса $\sim0{,}9$~Вт/м$^2$~\cite{loeb2012}. Дробление объективации на $O(1)$-акты ограничивает максимум одним битом, но не восстанавливает нулевой тариф. FIRAS~\cite{fixsen1996}, земной энергобаланс и лабораторный детальный баланс поэтому являются ограничениями на допустимое правило теплового регистра.

\paragraph{Спонтанная эмиссия и граница тарифа.}
Микроскопически тепловой фотон и фотон флуоресценции рождаются спонтанной эмиссией; одно название ансамбля не определяет тариф. Нужно задать подсистему акта, pointer-базис и физический локальный регистр подготовки (\Sref{sec:conditioning}). Контролируемая проверка $|{+}\rangle_{\mathrm{ptr}}$ в \Sref{sec:experiment} задаёт эти данные явно; для равновесной моды их онтическая спецификация остаётся условием модели.

\paragraph{Соотношение с существующими данными.}
Опубликованного \emph{прямого} экспериментального bound на тариф по $T$ и постулат $Q_i=kT\ln(1/p_i)$ нет; ниже~--- почему смежные измерения модель не исключают и не дублируют предложенный тест (\Sref{sec:experiment}).

\textbf{Оптика.} Однофотонная калориметрия на TES (Transition Edge Sensor; датчик на границе перехода)~\cite{karimi2020} разрешает энергию фотона 1550~нм до $\sim 0{,}1$--$1$~эВ, но при $T_{\mathrm{bath}}\sim 0{,}1$--$0{,}3$~K. При тарифе ``температура среды акта'' избыток $kT\ln 2$ при таком $T$ ($\sim 10^{-5}$~эВ на бит) на три--четыре порядка ниже разрешения TES. Холодные TES-данные \emph{не ограничивают} тариф по $T$; heralded-MZ тест (\Sref{sec:experiment}) с event-triggered калориметрией существующие однофотонные измерения не повторяет.

\textbf{Лазерная радиометрия.} Первичная шкала оптической мощности сравнивает тёплые тепловые детекторы с криогенными радиометрами электрического замещения ($T_{\mathrm{bath}}\sim 4$--$5$~K, тариф $kT\ln 2$ там $\sim 0{,}03\%$ от $E_{\gamma}$); согласие служб и типов детекторов~--- 0,02--1,2\%. Наивное чтение ``когерентный лазер $\Rightarrow$ лишние $\sim 2\%$ тепла при 300~K'' \emph{не} следует из модели: тариф платит не когерентность поля в базисе амплитуд, а суперпозицию \emph{в pointer-базисе акта} (\Sref{sec:pointer}) --- $|{+}\rangle_{\mathrm{ptr}}$ vs $I/2_{\mathrm{ptr}}$. Обычный измеритель мощности: один пучок, одно пятно, без делителя 50/50 и без контролируемой межзонной фазы; это не протокол $|{+}\rangle_{\mathrm{ptr}}$ из \Sref{sec:experiment} и не bound на тариф по $T$ для такого состояния. Предложенный heralded-MZ тест (\Sref{sec:experiment}) --- дифференциал $|{+}\rangle_{\mathrm{ptr}} - I/2_{\mathrm{ptr}}$ на однофотонных кубитах; радиометрия его не воспроизводила.

\textbf{Bounds на нагрев моделей коллапса.} Опыты по GRW/CSL/DP~\cite{bassi2013,donadi2021} (холодные атомы, кантилеверы, подземные установки) ограничивают \emph{любой} лишний нагрев от спонтанной локализации~--- ``фоновый'' нагрев самой модели коллапса. Heat from Bit спрашивает другое: сколько тепла даёт \emph{необратимая фиксация pointer-исхода}, если перед актом была когерентность в pointer-базисе ($H>0$). Существующие bounds опровергли бы Heat from Bit только если в тех установках при каждом событии гибло много pointer-битов. Обычно этого не происходит: система либо уже классическая смесь ($H\approx 0$), либо ``большая'' суперпозиция не в том pointer-базисе, где считается тариф~--- тогда $H=0$ и дополнительный нагрев по Heat from Bit тоже $\approx 0$. Измеренный нулевой нагрев в GRW/CSL/DP-тестах модель при тарифе по $T$ не запрещает. В \Sref{sec:experiment} $H$ и число битов за акт делают большими намеренно.

\paragraph{Сильный глобально обусловленный вариант.}
Если вместо рабочего локального правила принять вариант~(iii) из \Sref{sec:conditioning}, измерение Алисы удалённо выбирает для Боба ансамбль $\{|0\rangle,|1\rangle\}$ или $\{|+\rangle,|-\rangle\}$ при одной и той же $\rho_B=I/2$. Тогда идеализированный тепловой дифференциал Боба совпадает с MZ-сигналом~\eqref{eq:signal_power}; при параметрах \Sref{sec:experiment} время $\tau=(\mathrm{NEP}/\Delta P)^2\approx83$~с означает лишь $\mathrm{SNR}\approx1$, а не надёжно переданный бит. Такой результат был бы управляемой сверхсветовой сигнализацией и потребовал бы отдельного глобального закона причинного порядка, например выделенного слоения. Это \emph{не} следствие нелинейности $C_r$ и не часть основной версии модели. Бомовская механика показывает возможность выделенного слоения при скрытой нелокальности, но сама не даёт управляемой сигнализации; релятивистская flash-модель rGRWf, напротив, построена без предпочтительного слоения~\cite{bohm1952,tumulka2006}. Поэтому эти теории служат сравнением, а не доказательством необходимости выбранной конструкции.

\paragraph{Фальсификация.}
Прямая проверка модели~--- \Sref{sec:experiment}.

\paragraph{Сравнение с соседними подходами.}
От измерительных машин~\cite{elouard2017,elouard2018,yi2017,bresque2021,dassonneville2026,nogo2026} модель Heat from Bit отличается автономией акта без аппарата; от квантовой термодинамики проекций~\cite{kammerlander2016}~--- постулатом \emph{объективного} коллапса и тарифом в pointer-базисе акта, а не унитарным измерительным прибором; от информационных резервуаров~\cite{mandal2012,mandal2013,deffner2013,ptaszynski2019}~--- дуальностью (бит как носитель энергии при $\Delta S_{\mathrm{tot}}=0$); от GRW/CSL/DP~\cite{bassi2013}~--- организацией нагрева по потоку битов ($kT\ln 2$ на $\Delta N_{\mathrm{bit}}$): в стандартных формулировках их нагрев диффузен и не коррелирует с информационными актами (диффузный фон сам по себе модель Heat from Bit не отменяет).

\paragraph{Локальное первое начало.}
Первое начало нарушается лишь локально, в суб-$\tau_{\mathrm{th}}$ акте неунитарной фиксации (\Sref{sec:origin}); глобально второе начало сохраняется~--- $\delta Q$ рассеивается в резервуаре с энтропией $k\Delta H$. По отдельному событию дисбаланс энергии ниже разрешения калориметра; heralded MZ измеряет систематический избыток по ансамблю (\Sref{sec:experiment}).

\vspace{0.4em}
\begin{samepage}
\noindent\refstepcounter{table}\label{tab:comparison}%
{\centering
\textbf{Таблица~\thetable.} Сравнение стандартной КМ, моделей GRW/CSL/DP и модели Heat from Bit.\par\vspace{0.25em}
{\footnotesize
\setlength{\tabcolsep}{3pt}
\renewcommand{\arraystretch}{1.12}
\resizebox{\textwidth}{!}{%
\begin{tabular}{|>{\raggedright\arraybackslash}p{2.8cm}|c|c|c|}
\hline
 &
\makecell{\textbf{Стандартная}\\\textbf{КМ}} &
\makecell{\textbf{GRW/CSL/DP}} &
\makecell{\textbf{Heat from Bit}} \\
\hline
Объективный коллапс & нет & да & да \\
\hline
Первое начало & да & нет & \makecell{локально\\в акте} \\
\hline
Второе начало & да & да & да \\
\hline
Тепло при объективном акте & 0 & диффузный нагрев & $\langle Q\rangle=kTH$ \\
\hline
Отличительный тест & --- & \makecell{Gran Sasso;\\LISA Pathfinder;\\ультрахолодные атомы} & heralded MZ \\
\hline
\end{tabular}%
}
}
\par}
\end{samepage}
\vspace{0.6em}

\bibliographystyle{unsrturl}
\bibliography{refs}

\clearpage
\appendix
\renewcommand{\thesection}{\Alph{section}}
\section{Температура}\label{app:temp}

Тариф Heat from Bit использует температуру \emph{среды акта}~$T$~--- термодинамическую температуру резервуара, куда рассеивается платёж $\delta Q$ (термостат поглотителя, газ в ячейке, подложка микроострова). Во всех расчётах и экспериментах основного текста $T$~--- именно эта температура; $\mu_{\mathrm{bit}} = kT\ln 2$.

\paragraph{Альтернативные гипотезы.}
Помимо тарифа по температуре среды акта, модель допускает \emph{альтернативный} выбор параметра $T$ в $Q_i = kT\ln(1/p_i)$: «цена бита» не обязана совпадать с термостатом поглотителя. Ниже~--- две часто обсуждаемые версии; какую реализует природа, должен показать эксперимент.

\paragraph{Температура Унру/Риндлера.}
\begin{equation}
T_U = \frac{\hbar a}{2\pi c k_B}
\end{equation}
--- температура соответствующего горизонта ускоренного наблюдателя. Горизонт~--- не альтернативный носитель тепла в лаборатории: $T_U$ задаёт параметр тарифа, а не термостат болометра.

\paragraph{Температура де~Ситтера.}
\begin{equation}
T_{\mathrm{dS}} = \frac{\hbar H}{2\pi k_B}
\end{equation}
--- температура горизонта де~Ситтера ($H$~--- параметр Хаббла). Космологический вариант $T_{\Lambda} \approx 3{,}5$~K (прил.~\ref{app:xe}) относится к тому же классу альтернатив.

\paragraph{Геометрический дуал.}
На горизонте Риндлера/де~Ситтера та же бухгалтерия читается как прирост площади
\begin{equation}
\delta A = 4\ell_P^2\ln 2
\end{equation}
--- геометрический дуал платы $kT\ln 2$ (\Sref{sec:thermo}).

\paragraph{Калориметрия и различимость.}
Калориметр чувствует лишь тепло в лабораторной среде и потому напрямую проверяет только тариф по $T$ среды акта. $T_U$ и $T_{\mathrm{dS}}$ крайне малы: ожидаемое $\Delta Q$ лежит далеко ниже чувствительности, и нулевой сигнал калориметра с ними совместим~--- указание, что лабораторный тариф не реализован, а не провал модели. Различение $T$ vs.\ $T_U$ vs.\ $T_{\mathrm{dS}}$~--- косвенное (масштаб $\Delta Q$, космологические тесты).

\paragraph{Пример: $^{129}$Xe при $T_{\Lambda}$.}
Для SEOP-ячейки (прил.~\ref{app:xe}) при $T_{\Lambda} \approx 3{,}5$~K сигнал в $\sim 100$ раз слабее, чем при $T \approx 350$~K: $\Delta T \approx 0{,}1$~K, $\Delta P \approx 56$~Па, SNR $\sim 50$ vs.\ $\sim 10^3$.

\section{Калориметрия объективации ядерной когерентности ($^{129}$Xe)}\label{app:xe}

Прямая проверка тарифа инкассации на макроскопическом числе спинов, где $\langle\hat{H}\rangle$ и RF-мощность подавлены: SEOP-ячейка с гиперполяризованным $^{129}$Xe в поле $B_0 \sim 2$~мТл ($\hbar\omega_{\mathrm{Zeeman}} \sim 10^{-28}$~Дж на спин~--- на $\sim 7$ порядков ниже сигнала H1).

\paragraph{Pointer-базис и тариф.}
Zeeman-базис $|{\uparrow}\rangle,|{\downarrow}\rangle$~--- pointer-базис: градиент $B_0$ связывает энергию с $\hat\sigma_z$, среда декогерирует в этом базисе и размножает Zeeman-исходы (einselection). Каждое ядро~--- локализованная подсистема: фиксация Zeeman-исхода из $|{+}\rangle$ меняет $\rho$~--- как $|{+}\rangle_{\mathrm{ptr}}$ у поглотителя (\Sref{sec:experiment}); критерий \Sref{sec:pointer} выполнен поспиново; уничтожаемая когерентность в нём при поляризации $\varepsilon$:
\begin{equation}
H(\varepsilon) = \ln 2 - h_2\!\left(\tfrac{1+\varepsilon}{2}\right) \quad \text{(нат)}, \qquad H(0{,}5) \approx 0{,}131,
\end{equation}
где $h_2$~--- бинарная энтропия Шеннона. \textbf{H1} (Heat from Bit): $\langle Q\rangle = N_{\mathrm{Xe}}\, k T\, H(\varepsilon)$ при \emph{необратимой} гибели поперечной когерентности (среднее по $\sim 10^{21}$ актам); различие $Q_i$ по исходам здесь не различимо. \textbf{H0} (стандартная КМ + бесплатная унитарность): ноль везде сверх фонов. Тариф по температуре среды: $T \approx 350$~K (прил.~\ref{app:temp}).

\paragraph{Crusher: унитарность и pay-per-loss.}
Стандартный crusher-градиент~--- унитарный spoiler: после $\pi/2$ импульс $G$ на время $\tau$ накапливает фазу $\varphi(z)=\gamma G z\tau$ ($\gamma/2\pi \approx 11{,}8$~МГц/Т для $^{129}$Xe); оператор $U=\exp(-\mathrm{i}\gamma G\tau\sum_j z_j\hat S_j^z)$. Интегральный FID гаснет, но $M_\perp$ перекодирована в пространственную «спираль» с $k=\gamma G\tau$; gradient echo ($-\,G$ на то же $\tau$) восстанавливает сигнал~--- по \Sref{sec:thermo} тариф на момент импульса не активируется. Подлинная необратимость~--- диффузия: спин, набравший фазу в $z_1$, уплывает в $z_2$, и обратный градиент компенсирует не ту фазу. Амплитуда gradient echo
\begin{equation}
A(t)=e^{-Dk^2 t}, \qquad k=\gamma G\tau, \qquad \tau_{\mathrm{irr}}\approx (Dk^2)^{-1},
\end{equation}
--- прямой счётчик ещё живой поперечной когерентности; $\tau_{\mathrm{irr}}$ управляется $G\tau$ (от долей миллисекунды до секунд). Согласованно с порогом скремблирования~\cite{ponfilenok2026threshold}: платит не обращение известного $U$, а трансвычислительно необратимое смешение $\sim 10^{21}$ броуновских траекторий. Принцип \emph{pay-per-loss}: тепло следует за огибающей истинно необратимой потери, а не за RF-импульсом:
\begin{equation}
Q(t)=Q_{\mathrm{tot}}\bigl[1-e^{-Dk^2 t}\bigr], \qquad
\frac{\mathrm{d} Q}{\mathrm{d} t}=Dk^2 e^{-Dk^2 t}\,Q_{\mathrm{tot}}, \qquad
\frac{\mathrm{d} Q}{\mathrm{d} t}\;\parallel\;-\frac{\mathrm{d} A}{\mathrm{d} t}.
\end{equation}
Ноль на $\pi/2$ и на crusher-импульсе; подъём с задержкой $\sim\tau_{\mathrm{irr}}$. Отдельно $T_2$ (прочая декогеренция, секунды) и $T_1$ (распад поляризации, минуты~--- медленный фон).

\paragraph{Сигнал (базовая конфигурация).}
Ячейка 0{,}5~л, смесь 50\% Xe (обогащение $^{129}$Xe до 85\%) + 50\% N$_2$, 2~атм, $T = 350$~K, $\varepsilon = 0{,}5$. $N \approx 8{,}5 \times 10^{21}$, $D\approx 6\times 10^{-6}$~м$^2$/с (1~атм Xe). Тогда
\begin{equation}
Q_{\mathrm{tot}} = N\, k T\, H(\varepsilon) \approx 5{,}4~\text{Дж}.
\end{equation}
$C_v \approx 0{,}55$~Дж/К, $\Delta T_{\mathrm{tot}}\approx 10$~K, $\Delta P_{\mathrm{tot}}/P \approx 2{,}8\% \approx 5{,}6$~кПа (интеграл по полной необратимой потере). При мягком crusher'е ($G=10$~мТл/м, $\tau=5$~мс, $k\approx 3{,}7\times 10^3$~рад/м, $\tau_{\mathrm{irr}}\sim 0{,}5$--1~с) пиковая $\mathrm{d} Q/\mathrm{d} t\sim Q_{\mathrm{tot}}/\tau_{\mathrm{irr}}\sim 5$--10~Вт, а не киловаттный фронт на миллисекундах crusher'а.

\paragraph{Установка.}
Стандартная SEOP-ячейка с каплей Rb, поле $B_0 \sim 2$~мТл, накачка 795~нм. Добавить: дифференциальный датчик давления ($\lesssim 1$~Па, $\ge 1$~кГц); термопара; RF $\pi/2$ ($\sim 23$~кГц); crusher и gradient echo ($\pm G$); катушка приёма для $A(t)$; CPMG; синхронная запись RF/АЦП/$P(t)$/FID/$A(t)$.

\paragraph{Протоколы.}
\textbf{A} (pay-per-loss): накачка $\to$ $\pi/2$ + crusher ($G$, $\tau$) $\to$ запись $P(t)$ и $A(t)$ 60~с. H1: $Q(t)=Q_{\mathrm{tot}}[1-e^{-Dk^2 t}]$; $\mathrm{d} Q/\mathrm{d} t \parallel -\mathrm{d} A/\mathrm{d} t$; $\tau_{\mathrm{irr}}$ масштабируется как $(G\tau)^{-2}$ при фиксированном $Q_{\mathrm{tot}}$. \textbf{B} (контроль унитарности): (i)~$\pi/2 \to$ задержка $\sim 100$~мкс $\to$ $-\pi/2$; (ii)~$\pi/2 \to G\tau \to -G\tau \to -\pi/2$ (gradient echo); H0/H1: $P(t)$ и $A(t)$ плоски. \textbf{C} (дифференциал по $k$): два прогона с $G\tau$ и $3G\tau$ при равном RF/Joule-депозите; тривиальный нагрев сокращается, $\tau_{\mathrm{irr}}$ сдвигается в $\sim 9$ раз. \textbf{D} (CPMG): CPMG с интервалом $t_{\mathrm{cp}}$ подавляет диффузию как $t_{\mathrm{cp}}^2$; $\mathrm{d} Q/\mathrm{d} t \propto$ эффективной скорости потери эха; разность ``crusher''$-$``CPMG$\to$crusher'' изолирует необратимую компоненту. \textbf{Ось $\varepsilon$}: $Q_{\mathrm{tot}}\propto H(\varepsilon)\approx \varepsilon^2/2$.

\paragraph{Фоны.}
\noindent
{\small
\setlength{\tabcolsep}{4pt}
\begin{tabularx}{\textwidth}{|l|r|X|}
\hline
\textbf{Фон} & \textbf{Величина} & \textbf{Как давится} \\
\hline
Zeeman (2~мТл) & $\sim 10^{-7}$~Дж & подавлен по построению \\
RF-нагрев (23~кHz) & $< 10^{-6}$~Дж & пары A/B с равным RF \\
Joule градиентной катушки & $\propto G^2$, момент $\tau$ & не коррелирует с $A(t)$; C отменяет \\
Вихревые токи & Дж в стенках & момент импульса vs $\tau_{\mathrm{irr}}$ \\
Дрейф печки & $\pm$K за мин & слепая ячейка \\
\hline
\end{tabularx}
}

\paragraph{Чувствительность и исходы.}
Главный критерий~--- корреляция $P(t)\leftrightarrow A(t)$, а не одна ступень. При $\tau_{\mathrm{irr}}\sim 1$~с SNR по давлению ниже, чем у мгновенной ступени, но предсказание фальсифицируемо: задержанное тепло с $\tau_{\mathrm{irr}}\propto(G\tau)^{-2}$ не воспроизводится омическим нагревом катушки. \textbf{Подтверждение H1}: экспонента в A, ноль в B(i--ii), сдвиг $\tau_{\mathrm{irr}}$ в C, подавление нагрева в D при коротком $t_{\mathrm{cp}}$, скейлинг $\varepsilon^2$, тариф по $T$ (прил.~\ref{app:temp}). \textbf{Нулевой результат} $\Rightarrow$ bound на $T$; модель Heat from Bit ограничена снизу.
\paragraph{Соотношение с существующими данными (SEOP).}
Режим, близкий к базовой конфигурации (прил.~\ref{app:xe}), десятилетиями реализуется в SEOP и спин-фильтрах; по модели при $T \approx 350$~K и \emph{необратимой} диффузионной потере когерентности интеграл $\sim 5{,}4$~Дж за секунды означает $\Delta T\sim10$~K и мощность порядка $5$--$10$~Вт. Эффект такого масштаба трудно считать заведомо незаметным в ячейках, где температура и давление контролируются ради плотности паров Rb. Отсутствие опубликованной синхронной калориметрии не является численной границей, но многолетняя практика представляет существенное напряжение для тарифа по температуре среды. До ретроспективного анализа точности термометрии и режимов crusher/gradient echo приложение следует читать как stress test, а не как гарантированно открытое экспериментальное окно. Gradient echo и $\pi/2\to-\pi/2$ по постулату не платят (протокол~B); решающая проверка требует дифференциальной температуры/давления и синхронной записи $A(t)$ на существующей SEOP-ячейке.


\section{Fock-состояние и потолок $N_*$}\label{app:fock}
\paragraph{Граница тарифа: флуоресценция и bulk-поглощение.}
Микроскопически тепловой фотон и фотон флуоресценции рождаются спонтанной эмиссией; различие нельзя задавать свободным выбором ансамблевого разложения одной $\rho$. По правилу \Sref{sec:conditioning} тариф определяется подсистемой акта, pointer-базисом и физически реализованным локальным регистром подготовки. Контролируемые состояния~--- однофотонный $|{+}\rangle_{\mathrm{ptr}}$ (\Sref{sec:experiment}) или $|1\rangle_{\mathrm{Fock}}$ в заданной пространственной моде; для теплового источника онтический регистр должен быть указан отдельно. Наивное применение формулы ниже к каждому широкомодовому фотону давало бы до $kT\ln N_*\approx0{,}16$~эВ на bulk-поглощение, чего обычная фототермия не показывает. Дробление W-регистра по порогу скремблирования~\cite{ponfilenok2026threshold} ограничивает отдельный акт величиной $O(kT\ln2)$, но, как показано в \Sref{sec:discussion}, само по себе не восстанавливает детальный баланс. Поэтому геометрия ниже является намеренным тестом Fock-состояния, а перенос на флуоресценцию и равновесный свет остаётся открытым ограничением.


\paragraph{Соотношение с TES и радиометрией.}
Однофотонная калориметрия на TES (Transition Edge Sensor)~\cite{karimi2020} при $T_{\mathrm{bath}}\sim 0{,}1$--$0{,}3$~K разрешает энергию фотона 1550~нм до $\sim 0{,}1$--$1$~эВ; избыток $kT\ln 2$ ($\sim 10^{-5}$~эВ на бит) на три--четыре порядка ниже разрешения TES; на потолке Fock~--- $\lesssim 8{,}7\,kT\ln 2$. Холодные TES-данные \emph{не ограничивают} тариф по $T$; heralded-MZ тест (\Sref{sec:experiment}) с event-triggered калориметрией существующие однофотонные измерения не повторяет. Широкое когерентное пятно и потолок $kT\ln N_*$~--- отдельная постановка ($w\lesssim 0{,}2$~мкм); типовой ваттметр её не реализует. Для количественных оценок используется порог скремблирования~\cite{ponfilenok2026threshold}.

\paragraph{Pointer-ячейка и $r_{\mathrm{th}}$.}
$r_{\mathrm{th}}$~--- радиус пятна термализации; $r_{\mathrm{th}}^3$~--- объём элементарной pointer-ячейки, в которой фиксируется один исход локализации поглощения. Термализационный каскад за $t_{\mathrm{obj}}\approx10^{-12}$--$10^{-11}$~с набирает порядка $N_*$ независимых записей среды~\cite{ponfilenok2026threshold} и фиксирует $r_{\mathrm{th}}\approx1{,}5\times10^{-8}$--$2\times10^{-8}$~м. В heralded MZ (\Sref{sec:experiment}) две зоны задаются геометрией интерферометра ($N_{\mathrm{eff}}=2$); ниже $r_{\mathrm{th}}^3$ используется для подсчёта $N_{\mathrm{eff}}$ в делокализованной моде.

Проверка \emph{максимума} тарифа: делокализованное $|1\rangle_{\mathrm{Fock}}$ при субдифракционном фокусе рождает W-состояние $\sum_i c_i|e_i\rangle$ по $N_{\mathrm{eff}}$ pointer-ячейкам объёма $r_{\mathrm{th}}^3$. Объективация позиции уничтожает $\log_2 N_{\mathrm{eff}}$ бит:
\begin{equation}
\langle Q\rangle_{\mathrm{Fock}} = kT\ln N_{\mathrm{eff}} = kT\ln 2 \cdot \log_2 N_{\mathrm{eff}},
\qquad
N_{\mathrm{eff}} = V_{\mathrm{mode}}/r_{\mathrm{th}}^3.
\end{equation}
Бинарный heralded-MZ-тест (\Sref{sec:experiment}) даёт один бит ($N_{\mathrm{eff}}=2$); здесь~--- наращивание до потолка $N_*$.

\paragraph{Согласование с порогом скремблирования.}
По~\cite{ponfilenok2026threshold} при $C_c\sim10^{123}$ неструктурированный регистр упирается в потолок $N_*\sim\log_2 C_c\approx 408{,}6\approx409$. Which-path W-состояние на $N_{\mathrm{eff}}$ ячейках~--- как раз такой регистр: требуем
\begin{equation}
N_{\mathrm{eff}} \le N_* \approx 409
\qquad\Rightarrow\qquad
\log_2 N_{\mathrm{eff}} \le \log_2 N_* \approx 8{,}7.
\end{equation}
Иначе уже which-path-запутанность превысила бы порог до термализации. Отсюда жёсткий предел на объём моды: $V_{\mathrm{mode}}\le N_*\,r_{\mathrm{th}}^3$. При $r_{\mathrm{th}}\approx 15$--$20$~нм и толщине поглощения $\sim r_{\mathrm{th}}$ радиус оптического пятна $w\lesssim 0{,}2$~мкм; расфокус в микрон--миллиметр \emph{запрещён} как единый когерентный W-регистр (мода должна быть субдифракционно сжата или модово отфильтрована).

\paragraph{Оценка сигнала.}
Геральдированные $|1\rangle_{\mathrm{Fock}}$ (SPDC + SNSPD) при $N_{\mathrm{eff}}=N_*$: $\langle Q\rangle \le kT\ln N_* \approx 8{,}7\,kT\ln 2$. При $T=300$~K: $\lesssim 0{,}16$~эВ ($\sim 20\%$ от $E_{\gamma}$). При $T=77$~K: $\lesssim 0{,}04$~эВ ($\sim 5\%$). Для сравнения, один бит heralded MZ (\Sref{sec:experiment}): $\sim 2{,}2\%$ от $E_{\gamma}$.

\paragraph{Установка.}
\begin{itemize}
\item \textbf{Источник}: геральдированные $|1\rangle_{\mathrm{Fock}}$ на $\lambda = 1550$~нм; $R=10^{8}$~с$^{-1}$ (тот же класс источника, \Sref{sec:experiment}).
\item \textbf{Геометрия}: субдифракционный фокус $w\lesssim 0{,}2$~мкм на микроострове~--- $N_{\mathrm{eff}}\to N_*$; контрольные прогоны с раздуванием $w$ (ожидаемое насыщение).
\item \textbf{Поглотитель}: тот же класс, что в \Sref{sec:experiment}; предпочтителен $T\sim 300$~K для максимального $\Delta P/P$.
\item \textbf{Базовая линия}: чередование с heralded-MZ $|{+}\rangle_{\mathrm{ptr}}$ (\Sref{sec:experiment}) на том же пятне~--- проверка Fock$-$MZ: наращивание от $kT\ln 2$ до $kT\ln N_*$.
\end{itemize}

\paragraph{Что фиксируем.}
\begin{enumerate}
\item Избыток Fock $|1\rangle_{\mathrm{Fock}}$ над фоном $I/2_{\mathrm{ptr}}$ и над heralded-MZ $kT\ln 2$.
\item Fock$-$MZ при $w\lesssim 0{,}2$~мкм: $\langle Q\rangle$ растёт от одного бита до $\log_2 N_*$ бит.
\item \textbf{Насыщение по фокусу}: при раздувании $w$ избыток \emph{не} растёт пропорционально площади сверх $kT\ln N_*$~--- ключевой тест потолка $N_*$.
\item Линейность по $T$; материальный сдвиг через $r_{\mathrm{th}}$.
\end{enumerate}
\textbf{Исходы.} Подтверждение: избыток на упёртом в $N_*$ пятне; насыщение при расфокусе. Рост без насыщения~$\Rightarrow$ bound на $N_*$ или на модель; нуль при достаточной чувствительности~$\Rightarrow$ bound на $T$.

\section{Предел бит на акт коллапса}\label{app:bitcap}

В~\Sref{sec:origin} оценка~\eqref{eq:tau_col} связывает $\tau_{\mathrm{col}}$ с тарифным платежом через соотношение неопределённостей. Здесь~--- другая фундаментальная величина: квантовый предел скорости эволюции Марголуса--Левитина~\cite{margolus1998} и связанная с ним граница хаоса Maldacena--Shenker--Stanford (MSS)~\cite{maldacena2016mss}. Ниже~--- \emph{идея} (не постулат работы): перенос этой границы на скорость объективного коллапса даёт неравенства на~$\tau_{\mathrm{col}}$,~$\tau_{\mathrm{bit}}$ и на тариф~$H$, умещающийся в один акт.

Теорема Марголуса--Левитина ограничивает частоту ортогонализации состояния энергией над основным уровнем: $\nu_\perp\le 2E/(\pi\hbar)$. В тепловом резервуаре характерный масштаб $E\sim kT$ даёт скорость порядка $kT/\hbar$~--- ту же тепловую шкалу, что и $\tau_{\mathrm{th}}^{-1}$. Граница MSS уточняет ляпуновский рост вневременных корреляторов (OTOC) в тепловой квантовой системе:
\begin{equation}
\lambda \le \frac{2\pi kT}{\hbar}=\frac{2\pi}{\tau_{\mathrm{th}}}.
\label{eq:mss}
\end{equation}
Это интенсивный предел: он не растёт с числом частиц (прил.~\ref{app:bdpt}); быстрее тепловой хаос идти не может.

Предположим, что скорость \emph{неунитарной} фиксации pointer-исхода не превосходит эту границу:
\begin{equation}
\frac{1}{\tau_{\mathrm{col}}} \le \frac{2\pi kT}{\hbar},
\qquad
\tau_{\mathrm{col}} \ge \frac{\hbar}{2\pi kT}=\frac{\tau_{\mathrm{th}}}{2\pi}.
\label{eq:tau_col_mss}
\end{equation}
При $T=300$~K нижняя граница составляет $\hbar/(2\pi kT)\approx4\times10^{-15}$~с.

Тарифный платёж за один бит~--- $\delta Q=kT\ln 2$, энтропия $\Delta S=k\ln 2$. Отождествляя ``время бита'' с интервалом, на котором при скорости, не превышающей~$\lambda$, набирается информационная мера~$\ln 2$,
\begin{equation}
\tau_{\mathrm{bit}} \ge \frac{\ln 2}{\lambda_{\max}}=\frac{(\ln 2)\,\hbar}{2\pi kT}=\frac{\tau_{\mathrm{th}}\ln 2}{2\pi}.
\label{eq:tau_bit}
\end{equation}
При $T=300$~K: $\tau_{\mathrm{bit}}\ge2{,}8\times10^{-15}$~с. В отличие от верхней оценки~$\tau_{\mathrm{col}}$ из~\eqref{eq:tau_col} (окно допуска, в которое должен уместиться платёж),~\eqref{eq:tau_bit}~--- нижняя граница времени, сопряжённого максимальной скорости теплового хаоса и одному биту энтропии; $\mu_{\mathrm{bit}}=kT\ln 2$ остаётся тарифом энергии, а не новой динамикой.

Подчеркнём статус конструкции. (i)~MSS и Марголус--Левитин ограничивают \emph{унитарную} эволюцию и хаос; перенос~\eqref{eq:mss} на скорость \emph{объективного} коллапса~--- дополнительная гипотеза, не следствие этих теорем. (ii)~Равенство в~\eqref{eq:tau_col_mss}--\eqref{eq:tau_bit} достигалось бы лишь при насыщении границы; работа требует только неравенств. (iii)~Heat from Bit проверяется калориметрией $\langle Q\rangle=kT H$ (\Sref{sec:experiment}), а не прямым измерением~$\tau_{\mathrm{bit}}$; неравенства задают масштаб, на котором стоит искать микроскопическую динамику фиксации pointer-исхода в рамках предела Планка--Марголуса--Левитина / MSS.

\textbf{Окно для $\tau_{\mathrm{col}}$.} Собирая нижнюю границу MSS~\eqref{eq:tau_col_mss} и мандельштам--таммовскую оценку~\eqref{eq:tau_col} (она уже строже эвристики $\tau_{\mathrm{col}}\lesssim\tau_{\mathrm{th}}$ из начала~\Sref{sec:origin}, поскольку $1/(2\ln 2)<1$),
\begin{equation}
\frac{\tau_{\mathrm{th}}}{2\pi}
\le
\tau_{\mathrm{col}}
\lesssim
\frac{\tau_{\mathrm{th}}}{2\ln 2}.
\label{eq:tau_col_window}
\end{equation}
При $T=300$~K ($\tau_{\mathrm{th}}\approx2{,}5\times10^{-14}$~с):
\[
4\times10^{-15}~\text{с}
\le
\tau_{\mathrm{col}}
\lesssim
1{,}8\times10^{-14}~\text{с}.
\]

\textbf{Обобщение на тариф $H$.} Выше~--- случай одного бита ($H=\ln 2$). При объективном коллапсе с произвольным тарифом $H$ (\Sref{sec:thermo}; $H$ в натах) платёж $\delta Q=kT H$, и мандельштам--таммовское окно сжимается:
\begin{equation}
\tau_{\mathrm{col}} \lesssim \frac{\hbar}{2kT H}=\frac{\tau_{\mathrm{th}}}{2H}.
\label{eq:tau_col_mt_H}
\end{equation}
Параллельно, чтобы за время акта успеть ``протолкнуть'' тарифные $H$~нат когерентности при скорости, не превышающей~$\lambda_{\max}$, нужно
\begin{equation}
\tau_{\mathrm{col}} \ge \frac{H}{\lambda_{\max}}=\frac{H\,\hbar}{2\pi kT}=\frac{H}{\ln 2}\,\tau_{\mathrm{bit}}^{(1)},
\label{eq:tau_H}
\end{equation}
где $\tau_{\mathrm{bit}}^{(1)}$~--- нижняя граница~\eqref{eq:tau_bit} на один бит. Оценка~\eqref{eq:tau_col_mss} ($\tau_{\mathrm{col}}\ge\tau_{\mathrm{th}}/(2\pi)$)~--- ещё более слабый низ при $H\le 1$; при $H\ge 1$ доминирует~\eqref{eq:tau_H}. Собирая~\eqref{eq:tau_H} и~\eqref{eq:tau_col_mt_H},
\begin{equation}
\frac{H\,\tau_{\mathrm{th}}}{2\pi}
\le
\tau_{\mathrm{col}}
\lesssim
\frac{\tau_{\mathrm{th}}}{2H}.
\label{eq:tau_col_window_H}
\end{equation}
Окно непустое лишь при $H\le\sqrt{\pi}$ (нат), то есть $\lesssim 2{,}6$~бит. При $H=\ln 2<1$ нижняя граница~\eqref{eq:tau_H} слабее~\eqref{eq:tau_col_mss}, и восстанавливается~\eqref{eq:tau_col_window}; при $H\gtrsim\sqrt{\pi}$ (например потолок Fock $H\sim\ln N_*$ в прил.~\ref{app:fock}) верх уже ниже MSS-низа по~\eqref{eq:tau_H}: либо акт не может одновременно накопить весь~$H$ при~\eqref{eq:mss} и вместить платёж в $\Delta E\sim\hbar/(2\tau_{\mathrm{col}})$, либо фиксация дробится / выходит за тепловой шаг~$\tau_{\mathrm{th}}$.

\textbf{Физический смысл.} Граница $H\le\sqrt{\pi}$~--- не отдельный закон природы и не постулат Heat from Bit, а условие \emph{совместности} двух оценок для одного акта. Если коллапс (i)~не проталкивает больше $H$~нат когерентности быстрее теплового хаоса (MSS) и (ii)~весь тарифный платёж $kT H$ должен уместиться в допуск $\Delta E\sim\hbar/(2\tau_{\mathrm{col}})$, то за один такой акт нельзя уместить больше $\sqrt{\pi}$~нат ($\lesssim 2{,}6$~бит): иначе либо $\tau_{\mathrm{col}}$ слишком короток для накопления~$H$ при~\eqref{eq:mss}, либо $\Delta E$ акта меньше платежа. Это не запрет на большой~$H$ в модели вообще: при $H>\sqrt{\pi}$ ломается именно связка ``весь~$H$ за один MSS-ограниченный акт из~$\Delta E\sim\hbar/(2\tau)$''; обходы~--- более длинный акт, иной источник энергии платежа, дробление фиксации на шаги с меньшими~$H_i$.

\section{Время декогеренции при запутывании с тепловой средой (BDPT)}\label{app:bdpt}

Рассмотрим heralded-акт (\Sref{sec:experiment}): однофотонное path-состояние $|{+}\rangle_{\mathrm{ptr}}$ после поглощения в pointer-зоне. Сам фотон ($m=0$, $\mathrm{d}\tau=0$) BDPT-каналом не дефазируется~\cite{ponfilenok2026branch}; path-информация переносится на \emph{электрон} поглотителя (фотовозбуждение: валентный электрон или e--h пара), а не на целый атом сразу. Суперпозиция ``поглощено в зоне~0'' $|$ ``поглощено в зоне~1''~--- это суперпозиция электронных конфигураций в двух ячейках; начальная эффективная масса which-path-носителя
\begin{equation}
M_{\mathrm{eff}}^{(0)} \equiv m_0 \sim m_e
\quad\text{или}\quad
m_*\sim(2\text{--}5)\,m_e
\label{eq:m0}
\end{equation}
(масса e--h пары в полупроводнике). Ядерные координаты и соседние атомы подключаются на следующих шагах каскада.

В модели ветвезависимого собственного времени (BDPT)~\cite{ponfilenok2026branch} для which-path-ветвления с массой~$M$ ($E_{\mathrm{br}}=Mc^2$) потеря видности
\begin{equation}
\frac{V}{V_0}=e^{-\chi}, \qquad \chi=\Bigl(\frac{Mc^2}{\hbar}\Bigr)^{\!2}\delta t^2, \qquad \delta t^3=\xi\,t_P^2\,t,
\end{equation}
где $t_P=\sqrt{\hbar G/c^5}$, $\xi=\mathcal{O}(1)$; время декогеренции $t_{\mathrm{dec}}(M)$ определяется условием $\chi=1$:
\begin{equation}
t_{\mathrm{dec}}(M)=\frac{t_C(M)^3}{\xi\,t_P^2}, \qquad t_C(M)\equiv\frac{\hbar}{Mc^2}.
\label{eq:tdec_M}
\end{equation}
При $M=m_e$ (или $m_*$) $t_{\mathrm{dec}}\gg\tau_{\mathrm{th}}$: одного электронного приёмника недостаточно для BDPT-дефазировки за лабораторное окно.

\textbf{Критическая масса.} Условие $t_{\mathrm{dec}}(M)<\tau_{\mathrm{th}}$ (BDPT-шаг быстрее теплового) задаёт порог
\begin{equation}
M > M_{\mathrm{crit}} \equiv \frac{\hbar}{c^2\,\bigl(\xi\,t_P^2\,\tau_{\mathrm{th}}\bigr)^{1/3}}
= \Bigl(\frac{\hbar^3}{\xi\,t_P^2\,\tau_{\mathrm{th}}\,c^6}\Bigr)^{\!1/3},
\label{eq:Mcrit}
\end{equation}
полученный из~\eqref{eq:tdec_M} при $t_{\mathrm{dec}}(M_{\mathrm{crit}})=\tau_{\mathrm{th}}$. При $\xi=1$, $t_P=\sqrt{\hbar G/c^5}$ и $T=300$~K ($\tau_{\mathrm{th}}\approx2{,}5\times10^{-14}$~с): $M_{\mathrm{crit}}\approx2{,}8\times10^{-18}$~кг; при $T=77$~K: $\approx1{,}8\times10^{-18}$~кг; при $T=4$~K: $\approx6{,}6\times10^{-19}$~кг. Для произвольного $\xi$ все значения умножаются на $\xi^{-1/3}$. Так как $m_0\sim m_e\ll M_{\mathrm{crit}}$, на первом шаге BDPT не успевает; $M_{\mathrm{crit}}$ набирается каскадом~\eqref{eq:m_cascade}. В heralded MZ (\Sref{sec:experiment}) $M_{\mathrm{isl}}\sim10^{-18}$~кг при $T=300$~K находится в том же порядке величины, что $M_{\mathrm{crit}}$; при $M_{\mathrm{isl}}\gtrsim10^{-15}$~кг $t_{\mathrm{dec}}\ll\tau_{\mathrm{th}}$ уже при $\xi=\mathcal{O}(1)$.

\textbf{Каскад и наращивание массы.} Следующий ``приёмник'' pointer-информации (фонон, атом термализационного фронта) подключается за время порядка $\tau_{\mathrm{join}}\sim\tau_{\mathrm{th}}=\hbar/kT$. Пока $t<t_{\mathrm{dec}}(M_{\mathrm{eff}}^{(n)})$, which-path-суперпозиция ещё не BDPT-дефазирована; если $\tau_{\mathrm{join}}<t_{\mathrm{dec}}(M_{\mathrm{eff}}^{(n)})$, до подключения следующего приёмника масса успевает вырасти:
\begin{equation}
M_{\mathrm{eff}}^{(n+1)} = M_{\mathrm{eff}}^{(n)} + m_{\mathrm{rec}},
\qquad m_{\mathrm{rec}}\sim m_{\mathrm{at}},
\label{eq:m_cascade}
\end{equation}
где $m_{\mathrm{at}}$~--- масса атома/ионной ячейки поглотителя (Ge, InGaAs и т.п.). Каждый шаг укорачивает $t_{\mathrm{dec}}\propto M_{\mathrm{eff}}^{-3}$. Каскад продолжается, пока $M_{\mathrm{eff}}<M_{\mathrm{crit}}$ и $\tau_{\mathrm{join}}<t_{\mathrm{dec}}(M_{\mathrm{eff}}^{(n)})$; далее $t_{\mathrm{dec}}(M_{\mathrm{eff}})\lesssim\tau_{\mathrm{th}}$ или $M_{\mathrm{eff}}\to M_{\mathrm{isl}}$.

\textbf{Тепловая среда.} Параллельно действует \emph{средовой} канал~\cite{ponfilenok2026branch,hornberger2004tli}: запутывание с тепловыми модами резервуара на масштабе~$kT$ (газ, тепловое излучение, рассеяние) не входит в онтический канал $\delta t_a,\delta t_b$ BDPT, но рандомизирует фазу path-суперпозиции за
\begin{equation}
\tau_{\mathrm{dec}}^{(\mathrm{env})}\sim\tau_{\mathrm{th}}=\frac{\hbar}{kT}\approx 2{,}5\times10^{-14}~\text{с при }T=300~\text{K}.
\label{eq:tdec_env}
\end{equation}
При $T=77$~K $\tau_{\mathrm{dec}}^{(\mathrm{env})}\approx10^{-13}$~с; при $T=4$~K $\approx2\times10^{-12}$~с. Унитарное запутывание за $\tau_{\mathrm{dec}}^{(\mathrm{env})}$ и BDPT-дефазировка по нарастающему $M_{\mathrm{eff}}$ не активируют тариф (\Sref{sec:thermo}) и не суть объективного коллапса.

\textbf{Полное время декогеренции.} Различаем три масштаба. (i)~Мгновенное BDPT-время при фиксированной массе~$M$~--- формула~\eqref{eq:tdec_M}; это не полное время акта, а дефазировка при текущей~$M_{\mathrm{eff}}$. (ii)~Средовой канал~\eqref{eq:tdec_env}: фаза path-суперпозиции рандомизируется за один тепловой шаг~$\tau_{\mathrm{th}}$. (iii)~Время до BDPT-порога: каскад~\eqref{eq:m_cascade} должен набрать $M_{\mathrm{eff}}\gtrsim M_{\mathrm{crit}}$; наивно при $m_{\mathrm{rec}}\sim m_{\mathrm{at}}$ за каждый $\tau_{\mathrm{join}}\sim\tau_{\mathrm{th}}$
\begin{equation}
t_{\mathrm{cascade}}\sim\frac{M_{\mathrm{crit}}}{m_{\mathrm{at}}}\,\tau_{\mathrm{th}}\sim10^{-6}~\text{с}
\quad(T=300~\text{K},\; m_{\mathrm{at}}\sim m_{\mathrm{Ge}}),
\label{eq:t_cascade}
\end{equation}
после чего ещё~$t_{\mathrm{dec}}(M_{\mathrm{crit}})=\tau_{\mathrm{th}}$. В реальности термализационный фронт подключает много атомов параллельно; каскад до~$N_*$ записей среды занимает $t_{\mathrm{obj}}\sim10^{-12}$--$10^{-11}$~с (прил.~\ref{app:fock}). В heralded MZ при $T=300$~K средовой канал~\eqref{eq:tdec_env} быстрее BDPT-каскада, поэтому полное унитарное время декогеренции path-суперпозиции
\begin{equation}
t_{\mathrm{dec}}^{(\mathrm{tot})}\sim\tau_{\mathrm{dec}}^{(\mathrm{env})}\sim\tau_{\mathrm{th}}=\frac{\hbar}{kT}\approx2{,}5\times10^{-14}~\text{с}.
\label{eq:tdec_tot}
\end{equation}

\textbf{Глубина скремблирования против ширины записи.} Здесь два последовательных этапа, а не одно и то же. (i)~\emph{Запись} (дарвинизм~\cite{zurek2009}): каскад запутывает $N_{\mathrm{part}}$ приёмников; каждый хранит копию одного which-path-бита (состояние GHZ-типа). Ширина записи~--- $N_{\mathrm{part}}$; сложность распутывания на этом этапе линейна, $C_{\mathrm{copy}}\sim N_{\mathrm{part}}$, а не $2^{N_{\mathrm{part}}}$. Даже при $N_{\mathrm{part}}\sim2\times10^{7}$ (масса~$M_{\mathrm{crit}}$) имеем $C_{\mathrm{copy}}\ll C_c$: одной ширины для порога объективности~\cite{ponfilenok2026threshold} недостаточно. (ii)~\emph{Скремблирование}: после записи хаотическая фононная динамика (столкновения, ангармонизм) перемешивает копии; столкновение на частицу~--- раз в~$\sim\tau_{\mathrm{th}}$. Невосстановимое ядро растёт со временем:
\begin{equation}
N_{\mathrm{scr}}(t)\sim\min\Bigl(N_{\mathrm{part}}(t),\;\bigl(t/\tau_{\mathrm{th}}\bigr)^3\Bigr),
\qquad
C_{\mathrm{eff}}(t)\sim
\begin{cases}
N_{\mathrm{part}}, & N_{\mathrm{scr}}\ll N_{\mathrm{part}}\ \text{(ещё копии)},\\[2pt]
2^{\,N_{\mathrm{scr}}}, & \text{после перемешивания ядра}.
\end{cases}
\label{eq:Nscr}
\end{equation}
Диагностики перехода копий$\to$скремблинг~--- эхо Лошмидта и OTOC~\cite{ponfilenok2026threshold}: пока запись дарвинистская, эхо дёшево; после перемешивания ядра~--- нет. Скорость зависит от уже набранного $N_{\mathrm{part}}$ по-разному для разных величин. (a)~\emph{Интенсивная} (на степень свободы) ляпуновская шкала не растёт с~$N$: $\lambda\sim kT/\hbar$ (верхняя граница MSS $\lambda\le2\pi kT/\hbar$). (b)~\emph{Экстенсивная} частота столкновений во всём объёме записи~$\sim N_{\mathrm{part}}/\tau_{\mathrm{th}}$~--- линейна по числу уже запутанных частиц. (c)~Рост \emph{перемешанного ядра} ограничен геометрией фронта: в~(3+1)D при баллистическом разбегании корреляций $N_{\mathrm{scr}}(t)\sim\min\!\bigl(N_{\mathrm{part}},\,(t/\tau_{\mathrm{th}})^3\bigr)$, откуда
\begin{equation}
\frac{\mathrm{d}N_{\mathrm{scr}}}{\mathrm{d}t}\sim
\begin{cases}
\dfrac{3}{\tau_{\mathrm{th}}}\Bigl(\dfrac{t}{\tau_{\mathrm{th}}}\Bigr)^{\!2}
\sim\dfrac{3\,N_{\mathrm{scr}}^{2/3}}{\tau_{\mathrm{th}}}, & N_{\mathrm{scr}}<N_{\mathrm{part}},\\[8pt]
0, & N_{\mathrm{scr}}=N_{\mathrm{part}};
\end{cases}
\label{eq:scr_rate}
\end{equation}
скорость растёт с уже перемешанным объёмом ($\propto N_{\mathrm{scr}}^{2/3}$), пока фронт не охватит всю запись. (d)~Время \emph{полного} скремблирования уже записанных $N_{\mathrm{part}}$ частиц (быстрый скремблер / random circuit) $t_*\sim\tau_{\mathrm{th}}\ln N_{\mathrm{part}}$~--- логарифмически зависит от числа частиц; при $N_{\mathrm{part}}\sim2\times10^{7}$ и $T=300$~K это $t_*\sim4\times10^{-13}$~с. После перемешивания ядра $C_{\mathrm{eff}}\sim2^{N_{\mathrm{scr}}}$. При $t\sim\tau_{\mathrm{th}}$: $N_{\mathrm{scr}}\sim1$, ещё этап~(i). Порог $C_c\sim10^{123}$ требует этапа~(ii) с $N_{\mathrm{scr}}\sim\log_2 C_c\sim4\times10^{2}$; по~\eqref{eq:scr_rate} это $t\sim(4\times10^{2})^{1/3}\,\tau_{\mathrm{th}}\sim7\,\tau_{\mathrm{th}}\sim2\times10^{-13}$~с, а с запасом на неполноту фронта~--- $t_{\mathrm{obj}}\sim10^{-12}$--$10^{-11}$~с (прил.~\ref{app:fock}). К наивному $t_{\mathrm{cascade}}\sim10^{-6}$~с (набор~$M_{\mathrm{crit}}$) этап~(ii) давно пройден: срабатывает не ``$2^{N_{\mathrm{part}}}$ сразу'', а рост $N_{\mathrm{scr}}(t)$ до~$\sim4\times10^{2}$ при уже достаточном $N_{\mathrm{part}}$.

\textbf{Связь с $\tau_{\mathrm{col}}$.} Доминирует шаг~\eqref{eq:tdec_env} и/или каскад~\eqref{eq:m_cascade} при $M_{\mathrm{eff}}\to M_{\mathrm{isl}}$; неунитарная фиксация pointer-исхода должна завершиться за $\tau_{\mathrm{col}}\lesssim\tau_{\mathrm{th}}$~--- до подключения следующего ``приёмника'' (см. выше).

\end{document}
